\documentclass[10pt,a4paper]{article}
\usepackage[utf8]{inputenc}
\usepackage[T1]{fontenc}
\usepackage[french]{babel}
\usepackage[margin=1.2cm]{geometry}
\usepackage{enumitem}
\usepackage{microtype}
\usepackage[hidelinks]{hyperref}
\setlength{\parindent}{0pt}
\setlength{\parskip}{0pt}
\renewcommand{\baselinestretch}{1.00}
\pagestyle{empty}
\setlist[itemize]{leftmargin=*,nosep,topsep=1pt}
\newcommand{\cvsection}[1]{\par\vspace{3pt}\textbf{\large #1}\par\vspace{0.5pt}\rule{\linewidth}{0.35pt}\vspace{1pt}}

\begin{document}
\begin{center}
{\Large\textbf{LÉA DUBOIS}}\\[2pt]
\textbf{Ingénieure de recherche IA \& Bioinformatique | NLP, Python, Biologie | Paris}\\[2pt]
lea.dubois@example.test \;|\; +33 6 00 00 00 02 \;|\; \href{https://example.test/lea-dubois}{example.test/lea-dubois}
\end{center}

\cvsection{Résumé}
Ingénieure de recherche en bioinformatique spécialisée en NLP biomédical, machine learning et analyse de données omiques. Expérience dans la transformation de publications, données cliniques dé-identifiées et données transcriptomiques en modèles évalués, outils reproductibles et résultats interprétables pour des équipes scientifiques. Recherche un poste en IA appliquée aux sciences du vivant à partir d'octobre 2026.

\cvsection{Expérience}
\textbf{Institut BioNexus -- Ingénieure de recherche, NLP biomédical} \hfill 2024--2026
\begin{itemize}
\item Construction d'un pipeline Python d'ingestion et de normalisation de publications PubMed, avec annotation guidée par les ontologies MeSH et contrôle qualité des corpus.
\item Fine-tuning de BioBERT sous PyTorch pour reconnaître gènes, pathologies et traitements, avec macro-F1 de 0,88, validation croisée et analyse d'erreurs avec les biologistes.
\item Développement d'une API FastAPI de recherche sémantique et d'extraction d'entités, conteneurisée avec Docker et suivie par MLflow, DVC et GitHub Actions.
\item Cadrage des besoins, protocoles d'évaluation et restitution des résultats auprès de chercheurs, bioinformaticiens et équipes produit dans un environnement international.
\end{itemize}
\textbf{Genomic Systems Lab -- Bioinformaticienne, stage de recherche} \hfill 2023--2024
\begin{itemize}
\item Analyse de données single-cell RNA-seq : contrôle qualité, normalisation, réduction de dimension, clustering et identification de marqueurs avec Scanpy, Seurat et R.
\item Production de notebooks reproductibles, visualisations scientifiques et synthèses bibliographiques pour interpréter les sous-populations cellulaires observées.
\end{itemize}

\cvsection{Projets}
\textbf{Extraction de relations biomédicales} | Python, PyTorch, Transformers, BioBERT, spaCy
\begin{itemize}
\item Entraînement d'un pipeline NER et extraction de relations sur 4 200 résumés annotés, avec suivi des expériences, analyse d'erreurs par entité et interprétation des confusions.
\end{itemize}
\textbf{Assistant RAG de littérature scientifique} | Python, FastAPI, Qdrant, LangGraph, APIs LLM
\begin{itemize}
\item Développement d'un assistant sur des articles biomédicaux combinant recherche hybride, reranking, réponses sourcées et benchmark de questions conçu avec des utilisateurs scientifiques.
\end{itemize}
\textbf{Analyse single-cell de réponses immunitaires} | Python, Scanpy, R, Seurat, scVI
\begin{itemize}
\item Analyse reproductible de 18 000 cellules, de la qualité des données à l'annotation des clusters, avec comparaison de méthodes d'intégration et étude des voies biologiques différentielles.
\end{itemize}

\cvsection{Compétences}
\textbf{NLP \& GenAI} : PyTorch, Transformers, BioBERT, SciBERT, spaCy, NER, RAG, LangGraph\\
\textbf{Bioinformatique} : single-cell RNA-seq, Scanpy, Seurat, scVI, Biopython, BLAST, PubMed, MeSH\\
\textbf{Machine Learning \& statistiques} : scikit-learn, XGBoost, classification, validation croisée, SHAP, expérimentation\\
\textbf{Data \& développement} : Python, R, SQL, pandas, FastAPI, Docker, Git, MLflow, DVC, Linux

\cvsection{Formation}
\textbf{Université Paris-Saclay -- Master Bioinformatique \& Biostatistiques} \hfill 2022--2024\\
Génomique, biologie des systèmes, statistiques, machine learning, traitement automatique du langage.\\
\textbf{Université de Strasbourg -- Licence Sciences de la vie} \hfill 2019--2022\\
Biologie moléculaire, génétique, immunologie, programmation scientifique.

\cvsection{Publications \& communication scientifique}
Co-autrice d'un poster sur l'extraction d'entités biomédicales présenté à JOBIM 2025 | Animation de deux ateliers internes sur la reproductibilité des analyses et l'évaluation des modèles NLP.

\cvsection{Langues}
Français natif | Anglais C1, rédaction scientifique | Allemand B1

\cvsection{Centres d'intérêt}
Biologie marine | Médiation scientifique | Randonnée | Photographie microscopique
\end{document}
